\documentclass[10pt,a4paper]{amsart}
\usepackage{amsmath,amssymb,amsthm,graphicx}
% Retained minimal preamble for the Stacks Project chapter more-morphisms.tex (More on Morphisms).
% Verbatim excerpt of the chapter preamble: Xy-pic setup (native preamble lines 14-18) and the native
% theorem declarations with the mathematical macros (native preamble lines 154-217). Copied mechanically;
% nothing added, removed or redefined. The wrapper class and the amsmath/amssymb/amsthm/graphicx base
% packages are supplied by the assembler, so the class/externaldocument part of the native preamble is omitted.
% For commutative diagrams we use Xy-pic
\usepackage[all]{xy}

% We use 2cell for 2-commutative diagrams.
\xyoption{2cell}
% Theorem environments.
%
\theoremstyle{plain}
\newtheorem{theorem}[subsection]{Theorem}
\newtheorem{proposition}[subsection]{Proposition}
\newtheorem{lemma}[subsection]{Lemma}

\theoremstyle{definition}
\newtheorem{definition}[subsection]{Definition}
\newtheorem{example}[subsection]{Example}
\newtheorem{exercise}[subsection]{Exercise}
\newtheorem{situation}[subsection]{Situation}

\theoremstyle{remark}
\newtheorem{remark}[subsection]{Remark}
\newtheorem{remarks}[subsection]{Remarks}

\numberwithin{equation}{subsection}

% Macros
%
\def\lim{\mathop{\mathrm{lim}}\nolimits}
\def\colim{\mathop{\mathrm{colim}}\nolimits}
\def\Spec{\mathop{\mathrm{Spec}}}
\def\Hom{\mathop{\mathrm{Hom}}\nolimits}
\def\Ext{\mathop{\mathrm{Ext}}\nolimits}
\def\SheafHom{\mathop{\mathcal{H}\!\mathit{om}}\nolimits}
\def\SheafExt{\mathop{\mathcal{E}\!\mathit{xt}}\nolimits}
\def\Sch{\mathit{Sch}}
\def\Mor{\mathop{\mathrm{Mor}}\nolimits}
\def\Ob{\mathop{\mathrm{Ob}}\nolimits}
\def\Sh{\mathop{\mathit{Sh}}\nolimits}
\def\NL{\mathop{N\!L}\nolimits}
\def\CH{\mathop{\mathrm{CH}}\nolimits}
\def\proetale{{pro\text{-}\acute{e}tale}}
\def\etale{{\acute{e}tale}}
\def\QCoh{\mathit{QCoh}}
\def\Ker{\mathop{\mathrm{Ker}}}
\def\Im{\mathop{\mathrm{Im}}}
\def\Coker{\mathop{\mathrm{Coker}}}
\def\Coim{\mathop{\mathrm{Coim}}}

% Boxtimes
%
\DeclareMathSymbol{\boxtimes}{\mathbin}{AMSa}{"02}

%
% Macros for moduli stacks/spaces
%
\def\QCohstack{\mathcal{QC}\!\mathit{oh}}
\def\Cohstack{\mathcal{C}\!\mathit{oh}}
\def\Spacesstack{\mathcal{S}\!\mathit{paces}}
\def\Quotfunctor{\mathrm{Quot}}
\def\Hilbfunctor{\mathrm{Hilb}}
\def\Curvesstack{\mathcal{C}\!\mathit{urves}}
\def\Polarizedstack{\mathcal{P}\!\mathit{olarized}}
\def\Complexesstack{\mathcal{C}\!\mathit{omplexes}}
% \Pic is the operator that assigns to X its picard group, usage \Pic(X)
% \Picardstack_{X/B} denotes the Picard stack of X over B
% \Picardfunctor_{X/B} denotes the Picard functor of X over B
\def\Pic{\mathop{\mathrm{Pic}}\nolimits}
\def\Picardstack{\mathcal{P}\!\mathit{ic}}
\def\Picardfunctor{\mathrm{Pic}}
\def\Deformationcategory{\mathcal{D}\!\mathit{ef}}
\newenvironment{reference}{}{}
\newenvironment{slogan}{}{}
\newenvironment{history}{}{}
\usepackage{cleveref}
\crefname{originallemma}{original lemma}{original lemmas}
\Crefname{originallemma}{Original lemma}{Original lemmas}
\makeatletter
\newlabel{morphisms-lemma-characterize-quasi-affine}{{01SM}{source}{}{sourcepointer-1}{}}
\newlabel{morphisms-lemma-characterize-quasi-affine@cref}{{[originallemma][0][]01SM}{[1][0][]source}{}{}{}}
\makeatother
\crefname{originallemma}{original lemma}{original lemmas}
\Crefname{originallemma}{Original lemma}{Original lemmas}
\makeatletter
\newlabel{morphisms-lemma-quasi-affine-permanence}{{054G}{source}{}{sourcepointer-2}{}}
\newlabel{morphisms-lemma-quasi-affine-permanence@cref}{{[originallemma][0][]054G}{[1][0][]source}{}{}{}}
\makeatother
\crefname{originallemma}{original lemma}{original lemmas}
\Crefname{originallemma}{Original lemma}{Original lemmas}
\makeatletter
\newlabel{schemes-lemma-compose-after-separated}{{01KV}{source}{}{sourcepointer-3}{}}
\newlabel{schemes-lemma-compose-after-separated@cref}{{[originallemma][0][]01KV}{[1][0][]source}{}{}{}}
\makeatother
\crefname{originallemma}{original lemma}{original lemmas}
\Crefname{originallemma}{Original lemma}{Original lemmas}
\makeatletter
\newlabel{schemes-lemma-quasi-compact-permanence}{{03GI}{source}{}{sourcepointer-4}{}}
\newlabel{schemes-lemma-quasi-compact-permanence@cref}{{[originallemma][0][]03GI}{[1][0][]source}{}{}{}}
\makeatother
\title{Four testing criteria for ind-quasi-affine morphisms and Zariski locality on the base}
\author{The Stacks Project Authors}
\date{stacks-project:more-morphisms.tex:local-snapshot-20261011; GNU FDL 1.2 or later}
\begin{document}
\maketitle
\noindent\textit{Editorial note.} Statement span is lines 19537-19557 of sources/stacks-more-morphisms/pinned-original-source.tex; proof span is lines 19559-19582, the complete author proof (22 body lines).\par\medskip
\noindent\textit{Editorial note.} Navigation stubs: 4 cited-but-not-included labels are supplied as non-load-bearing pointers with their pinned tags where the snapshot tag index has them.\par\medskip
\noindent\textit{Editorial note.} extra\_packages is empty on purpose: the retained preamble already loads xy with options and a second declaration clashes in the assembler wrapper.\par\medskip
\noindent\textit{Editorial note.} Text is retained verbatim, including the author's own typos and placeholders.\par\medskip
\noindent\textit{Editorial note.} Extends the lane's fpqc-local companion card L16080 (ind-quasi-affine-fpqc-local): this is the Zariski-side structural dictionary used by that descent proof; the in-particular locality sentence is retained.\par\medskip
\noindent\textit{Allowed standard prerequisites (editorial).}\begin{itemize}
\item Explicit prerequisite boundary. Definition of ind-quasi-affine morphism (chapter definition section), Morphisms, lemmas characterize-quasi-affine and quasi-affine-permanence; Schemes, lemmas quasi-compact-permanence and compose-after-separated.
\end{itemize}\medskip
\ifdefined\hypertarget\hypertarget{sourcepointer-1}{}\fi
\noindent\textit{Source reference.} Original lemma 01SM: sources/stacks-more-morphisms/pinned-original-source.tex, source lines 19582--19582. Cited by the retained statement or proof but not part of the collected spans; supplied only as a navigation pointer. The argument does not depend on its content.\par\medskip
\ifdefined\hypertarget\hypertarget{sourcepointer-2}{}\fi
\noindent\textit{Source reference.} Original lemma 054G: sources/stacks-more-morphisms/pinned-original-source.tex, source lines 19582--19582. Cited by the retained statement or proof but not part of the collected spans; supplied only as a navigation pointer. The argument does not depend on its content.\par\medskip
\ifdefined\hypertarget\hypertarget{sourcepointer-3}{}\fi
\noindent\textit{Source reference.} Original lemma 01KV: sources/stacks-more-morphisms/pinned-original-source.tex, source lines 19582--19582. Cited by the retained statement or proof but not part of the collected spans; supplied only as a navigation pointer. The argument does not depend on its content.\par\medskip
\ifdefined\hypertarget\hypertarget{sourcepointer-4}{}\fi
\noindent\textit{Source reference.} Original lemma 03GI: sources/stacks-more-morphisms/pinned-original-source.tex, source lines 19582--19582. Cited by the retained statement or proof but not part of the collected spans; supplied only as a navigation pointer. The argument does not depend on its content.\par\medskip
\begin{lemma}
\label{lemma-ind-quasi-affine-alternative-definition}
For a morphism of schemes $f : X \to Y$, the following are equivalent:
\begin{enumerate}
\item $f$ is ind-quasi-affine,
\item for every affine open subscheme $V \subset Y$ and
every quasi-compact open subscheme $U \subset f^{-1}(V)$,
the induced morphism $U \to V$ is quasi-affine. 
\item
for some cover $\{ V_j \}_{j \in J}$ of $Y$ by
quasi-compact and quasi-separated open subschemes
$V_j \subset Y$, every $j \in J$, and every quasi-compact
open subscheme $U \subset f^{-1}(V_j)$, the induced morphism
$U \to V_j$ is quasi-affine.
\item for every quasi-compact and quasi-separated open subscheme
$V \subset Y$ and every quasi-compact open subscheme
$U \subset f^{-1}(V)$, the induced morphism $U \to V$ is quasi-affine.
\end{enumerate}
In particular, the property of being an ind-quasi-affine morphism
is Zariski local on the base.
\end{lemma}

\begin{proof}
The equivalence (1) $\Leftrightarrow$ (2)
follows from the definitions and
Morphisms, Lemma \ref{morphisms-lemma-characterize-quasi-affine}.
For (2) $\Rightarrow$ (4), let $U$ and $V$ be as in (4). By
Schemes, Lemma \ref{schemes-lemma-quasi-compact-permanence}, the
induced morphism $U \to V$ is quasi-compact. Thus, for every affine
open $V' \subset V$, the fiber product $V' \times_V U$ is quasi-compact,
so, by (2), the induced map $V' \times_V U \to V'$ is quasi-affine.
Thus, $U \to V$ is also quasi-affine by
Morphisms, Lemma \ref{morphisms-lemma-characterize-quasi-affine}.
This argument also gives (3) $\Rightarrow$ (4): indeed, keeping the
same notation, those affine opens $V' \subset V$ that lie in one
of the $V_j$ cover $V$, so one needs to argue that the
quasi-compact map $V' \times_V U \to V'$ is quasi-affine.
However, by (3), the composition $V' \times_V U \to V' \to V_j$
is quasi-affine and, by
Schemes, Lemma \ref{schemes-lemma-compose-after-separated}, the map
$V' \to V_j$ is quasi-separated. Thus, $V' \times_V U \to V'$
is quasi-affine by
Morphisms, Lemma \ref{morphisms-lemma-quasi-affine-permanence}.
The final implications (4) $\Rightarrow$ (2) and (4) $\Rightarrow$ (3)
are evident.
\end{proof}

\end{document}
